\section{Conditional burst means}
\label{dist:sec:cond-means}

Three different averages of the burst size appear in this chapter. They differ
by factors of nearly two and a half, and they are easy to confuse, because in
ordinary usage all three are written $\E{\Kb\mid\cdot}$. This section
separates them and quantifies the difference; the results are collected in
\cref{dist:tab:threemeans} and plotted in \cref{dist:fig:threemeans}, and
\ChPop{} draws on all three.

\subsection{The overall expected burst size}
\label{dist:sec:overall-mean}

Whether a catastrophe has occurred by $t$ is recorded by the truth of
$\Wt>0$, the minimum burst size being $1$; that event has probability
$1-I(t)$. Decomposing $\E{\Wt}$ on it, and using
$\E{\Wt\mid\Wt=0}=0$,
\begin{equation}
\E{\Wt\mid\Wt > 0} = \frac{V(t)}{1-I(t)}.
\label{dist:eq:voveri}
\end{equation}
This is the mean burst size given that a burst has happened by $t$. The
expected size of an arbitrarily generated burst is its late-time value,
\begin{equation}
\E{\Kb}=\lim_{t\to\infty}\frac{V(t)}{1-I(t)}
=\frac{V_\infty}{1-b}
=\frac{\lambda-\mu}{\delta}+\frac{1}{1-b}
=\frac{a}{a-1},
\label{dist:eq:meank}
\end{equation}
the last two forms from \cref{dist:eq:vinf}. The distinction from $V_\infty$
itself matters: $V_\infty$ averages in the realisations that cleared
internally and released nothing, so it is \emph{not} the mean burst size
unless $\mu=0$.

\subsection{The apparent circularity}
\label{dist:sec:circularity}

What is the expected burst size of a cell that has \emph{not yet} burst?
Decomposing $\E{\Kb}$ on the event $\{\tau>t\}$,
\begin{align}
\E{\Kb} &= \E{\Kb\mid \tau > t}\pr{\tau > t}
+ \E{\Kb\mid \tau < t}\pr{\tau<t}\nonumber\\
&= \E{\Kb\mid\tau > t}\, I(t) + \frac{V(t)}{1 - I(t)}\lrb{1 - I(t)},
\label{dist:eq:kdecomp}
\end{align}
so that
\begin{equation}
\E{\Kb\mid \tau > t} = \frac{\E{\Kb} - V(t)}{I(t)}
=\frac{V_\infty/(1-b) - V(t)}{I(t)}.
\label{dist:eq:kcond}
\end{equation}
Letting $t\to\infty$ in \cref{dist:eq:kcond} appears to give
\[
\lim_{t\to\infty}\E{\Kb\mid \tau > t}
=\frac{V_\infty/(1-b)-V_\infty}{b},
\]
which after a little algebra returns $\E{\Kb}$ itself. That cannot be right.
Conditioning on $\tau>t$ as $t\to\infty$ selects the realisations that never
burst at all, and their burst size is zero, not typical.

The resolution is that \cref{dist:eq:kcond} mixes two conditionings. Its
numerator $\E{\Kb}-V(t)$ is the contribution to the burst size from
realisations that burst \emph{after} $t$; its denominator $I(t)$ is the
probability of not having burst by $t$, which includes the mass $b$ of
realisations that never will and whose burst size is zero. Written
consistently,
\begin{equation}
\E{\Kb\,\ind\{\text{burst}\}\mid\tau>t}
=\frac{V_\infty-V(t)}{I(t)}
\xrightarrow[t\to\infty]{}\frac{V_\infty-V_\infty}{b}=0,
\label{dist:eq:consistent}
\end{equation}
which is correct and unmysterious. The genuinely interesting limit conditions
on eventual burst as well:
\begin{equation}
\E{\Kb\mid \tau>t,\ \text{burst eventually}}
=\frac{V_\infty-V(t)}{I(t)-b},
\label{dist:eq:patient}
\end{equation}
of the form $0/0$ as $t\to\infty$. By l'H\^opital's rule, with $V'=\delta K$
and $I'=-\delta J$,
\begin{equation}
\lim_{t\to\infty}\E{\Kb\mid \tau>t,\ \text{burst eventually}}
=\lim_{t\to\infty}\frac{K(t)}{J(t)}
=\frac{\langle X^2\rangle_{\rm QS}}{\langle X\rangle_{\rm QS}}
=1+\frac{2\lambda}{\delta}(1-b),
\label{dist:eq:latemean}
\end{equation}
in agreement with \cref{dist:eq:latelimit}: the eventual burst size of a very
patient cell is the size-biased quasi-stationary mean.

No interchange of limits is needed to reach it. The defect cancels in the
ratio, so $K/J$ is at every time the ratio of the second conditional moment to
the first, and by \cref{dist:eq:condgeom} that conditional law is geometric
with ratio $P(t)$; hence
\begin{equation}
\frac{K(t)}{J(t)}=\frac{1+P(t)}{1-P(t)},
\label{dist:eq:kjP}
\end{equation}
an identity holding at every $t$, not only in the limit.
\Cref{dist:eq:latemean} is the value of \cref{dist:eq:kjP} at the endpoint
$P=1/a$, and the only limit taken anywhere in the argument is
$P(t)\uparrow1/a$ from \cref{dist:thm:geometric}.

\begin{table}[tbp]
\centering
\renewcommand{\arraystretch}{1.35}
\small
\begin{tabular}{@{}llc@{}}
\toprule
Quantity & Definition & At $(1,0.2,0.05)$ \\
\midrule
$V_\infty$ & $\E{\Kb\ind\{\text{burst}\}}$; includes non-bursters at $0$ & $13.99$ \\
$\E{\Kb\mid\text{burst}}$ & conditioned on eventual rupture; $=a/(a-1)$ & $17.23$ \\
$\lim_t\E{\Kb\mid\tau>t,\,\text{burst}}$ & size-biased quasi-stationary mean; $=(a{+}1)/(a{-}1)$ & $33.46$ \\
\bottomrule
\end{tabular}
\caption{Three means, three conditionings. The first two differ by the factor
$1-b$, the probability of eventual rupture; the second and third by the
size-biasing factor $(a+1)/a$ of \cref{dist:sec:size-bias}.}
\label{dist:tab:threemeans}
\end{table}

\begin{remark}[Where the circularity was]
\label{dist:rem:circularity}
The circularity was never in the mathematics; it was in the conditioning.
Each of the three quantities in \cref{dist:tab:threemeans} is a correct
answer to a different question, and only the middle one is what is usually
meant by ``the mean burst size''. An assay that observes a cell because it
ruptured measures the second; an assay that observes cells still infected at
a late time and waits measures the third.
\end{remark}

Three means, and with them the cell founded by one particle is completely
described. The two sections that follow ask how much of that survives when the
founder is not alone, and when the cell is not.

\begin{figure}[tbp]
\centering
\includegraphics[width=0.72\textwidth]{fig11_three_burst_means/N4a_4_three_burst_means}
\caption{Three distinct averages at $(\lambda,\mu,\delta)=(1,0.2,0.05)$,
drawn as trajectories. The unconditional mean release $V(t)$ approaches
$V_\infty=13.99$ and includes cells that never burst; the burst mean
$V_\infty/(1-b)=a/(a-1)=17.23$ is the horizontal line; conditioning on a
very late eventual burst approaches the size-biased limit
$(a+1)/(a-1)=33.46$. Keeping the conditioning explicit removes the apparent
circularity.}
\label{dist:fig:threemeans}
\end{figure}
